\chapter{Background}

\section{The \TransHyp{} Problem}

\subsection{Properties of Transversal Hypergraphs}

This section explores the sufficient and necessary conditions for a hypergraph to be the transversal hypergraph of another. The following results are based on \cite{Ber87}.

\begin{lemma}\label{lem:partition}
Let $G$ and $H$ be two Sperner hypergraphs on the same vertex set $V$. Then $H=\Tr(G)$ if and only if, for all partitions $(A,B)$ of $V$, where $A\cup B=V$ and $A\cap B=\emptyset$, the following conditions hold:
\begin{enumerate}
\item There exists either an $e_i\in E(H)$ such that $e_i\subseteq A$ or an $e_j\in E(G)$ such that $e_j\subseteq B$.
\item These two cases cannot happen simultaneously.
\end{enumerate}
\end{lemma}

This lemma establishes a fundamental property for recognizing transversal hypergraphs, which leads directly to the following corollaries.

\begin{corollary}
Let $G$ and $H$ be two Sperner hypergraphs on the same vertex set $V$. Then $H=\Tr(G)$ if and only if $G=\Tr(H)$.
\end{corollary}

\begin{corollary}
Let $H$ be a Sperner hypergraph. Then $\Tr(\Tr(H))=H$.
\end{corollary}

\subsection{Problem Representation}

We define \TransHyp{} as the problem of determining if two given hypergraphs are transversal to each other.

\begin{problem}[\textit{\texttt{Trans-Hyp}}]
\textbf{\textit{Input:}} Two hypergraphs $G=(V,E(G))$ and $H=(V,E(H))$.\\
\textbf{\textit{Output:}} Determine whether $G$ and $H$ are transversal to each other.
\end{problem}

Clearly, \TransHyp{} for two hypergraphs $G$ and $H$ is in co-NP, as shown in \cite{FK96}, since if $H$ is not the transversal hypergraph of $G$, this can be verified by a set of vertices $S\subseteq V$ such that:
\begin{equation}\label{eq:witness}
S\cap e_i\neq\emptyset\text{ for all }e_i\in E(G),\quad\text{and}\quad S\nsupseteq e_j\text{ for all }e_j\in E(H).
\end{equation}
Intuitively, $S$ is a transversal of $G$ (not necessarily minimal) that does not fully contain any hyperedge of $H$. Such a set $S$ acts as a \emph{certificate} for the non-transversality of the pair $(G,H)$, as discussed in \cite{Rau11}, although this perspective is inspired from \cite{FK96} and earlier papers. We note that this condition is symmetric: if $S\subseteq V$ satisfies the above for $(G,H)$, then $V\setminus S$ satisfies it for $(H,G)$.

Hypergraphs are one medium for studying the \TransHyp{} problem. Monotone Boolean functions (MBFs) offer an alternate representation of the same problem. A Sperner hypergraph corresponds to the disjunctive normal form (DNF) of an MBF, where each hyperedge represents a minimal true setting of the function. Similarly, the transversal hypergraph corresponds to the conjunctive normal form (CNF) of the dual MBF, with each clause corresponding to a minimal condition for satisfaction. Using the MBF language, we can define the decision problem as the following:

\begin{problem}[\textit{\texttt{DUAL}}]
\textbf{\textit{Input:}} Two monotone Boolean functions $f$ and $g$, in CNF forms.\\
\textbf{\textit{Output:}} Are $f$ and $g$ dual?
\end{problem}

\section{Fredman-Khachiyan Algorithm \FKA}

The \FKA{} algorithm checks whether two hypergraphs $G$ and $H$ are a transversal pair through a recursive decomposition that generates a \emph{recursion tree}, where each node represents a pair of hypergraphs being checked. At the root of the tree are the input hypergraphs $G$ and $H$. The algorithm proceeds by splitting the input into two smaller subproblems whenever necessary, with the edges of the tree connecting parent nodes to their child subproblems.

The algorithm begins by performing a Sperner reduction on $G$ and $H$, removing edges that are not minimal. After this, it checks the necessary conditions that all transversal pairs must satisfy as a base step, and they can be verified in polynomial time. If any of these preconditions fail, the algorithm immediately returns a witness based on \ref{lem:partition}. If the hypergraphs pass the precondition checks, the algorithm handles trivial instances of input hypergraphs (hypergraphs with a single edge containing 1 vertex or the empty hypergraph). In the recursive case, the algorithm selects the most frequent vertex in $G$ or $H$ as a pivot. Using this vertex, it partitions the edge sets of $G$ and $H$ into two smaller subproblems.

Importantly, if $G$ and $H$ are already transversal, these precondition checks do nothing to reduce computation and the algorithm decomposes further into smaller subproblems, creating child nodes in the recursion tree. The performance of \FKA{} is determined by the number of subproblems it generates, and this is reflected in the number of nodes of the recursion tree \FKA{} generates.

\cite{FK96-missing}

\subsection{Base Conditions for \FKA}

Let $G$ and $H$ be Sperner hypergraphs that are transversal to one another. \cite{DNU21}, \cite{FK96} outlines some fundamental properties that must hold.

\begin{proposition}\label{prop:intersect}
For every $e_i\in E(G)$ and $e_j\in E(H)$, $e_i\cap e_j\neq\emptyset$.
\end{proposition}

\begin{proof}
Assume that there exists an edge $e_i\in E(G)$ and $e_j\in E(H)$ such that $e_i\cap e_j=\emptyset$. In this case, $e_i$ does not intersect $e_j$, meaning that $e_i$ fails to act as a minimal hitting set for $H$, which contradicts the assumption that $G$ and $H$ form a transversal pair.
\end{proof}

\begin{proposition}\label{prop:support}
$\bigcup\{e_i:e_i\in E(G)\}=\bigcup\{e_j:e_j\in E(H)\}$.
\end{proposition}

\begin{proof}
Suppose $\bigcup\{e_i:e_i\in E(G)\}\neq\bigcup\{e_j:e_j\in E(H)\}$. This implies either
\[
\bigcup\{e_i:e_i\in E(G)\}\nsubseteq\bigcup\{e_j:e_j\in E(H)\}\quad\text{or}\quad\bigcup\{e_j:e_j\in E(H)\}\nsubseteq\bigcup\{e_i:e_i\in E(G)\}.
\]

If $\bigcup\{e_i:e_i\in E(G)\}\nsubseteq\bigcup\{e_j:e_j\in E(H)\}$, there exists some vertex $v\in\bigcup\{e_i:e_i\in E(G)\}$ such that $v\notin\bigcup\{e_j:e_j\in E(H)\}$. Since $v$ is not contained in any edge of $H$, removing $v$ from $H$ does not intersect any edge in $H$. Thus, $H$ is not affected by the removal of $v$, and $G$ and $H$ cannot form a transversal pair.

Conversely, if $\bigcup\{e_j:e_j\in E(H)\}\nsubseteq\bigcup\{e_i:e_i\in E(G)\}$, there exists some vertex $v\in\bigcup\{e_j:e_j\in E(H)\}$ such that $v\notin\bigcup\{e_i:e_i\in E(G)\}$. Let $e_j\in E(H)$ be an edge containing $v$. Since $v$ does not appear in any edge of $G$, removing $v$ from $G$ does not intersect any edge in $G$. Therefore, $G$ and $H$ cannot form a transversal pair.
\end{proof}

\begin{proposition}\label{prop:size}
$\max\{|e_i|:e_i\in E(G)\}\leq|H|$ and $\max\{|e_j|:e_j\in E(H)\}\leq|G|$.
\end{proposition}

\begin{proof}
First, we show that $\max\{|e_i|:e_i\in E(G)\}\leq|H|$. Let's assume that there exists an edge $e_i\in E(G)$ such that $|e_i|>|H|$. Let $e_i'$ be a proper subset of $e_i$ where $e_i'\cap e_j\neq\emptyset$ for every $e_j\in H$. Such a subset must exist. Suppose we remove $e_i'$ from $G$ and also remove every vertex in $e_i'$ from $H$. Then $e_i\neq\emptyset$, so $e_i$ still contains at least one vertex $v$. However, since all edges in $H$ are intersected, this implies that $e_i'$ is a transversal of $H$. As $G$ is Sperner, this would mean $e_i'=e_i$, which is a contradiction because $e_i$ was assumed to have vertices not in $e_i'$. Hence, such an $e_i$ cannot exist.

Now we show that $\max\{|e_j|:e_j\in E(H)\}\leq|G|$. Assume that there exists an edge $e_j\in E(H)$ such that $|e_j|>|G|$. Let $e_j'$ be a proper subset of $e_j$ such that $e_j'\cap e_i\neq\emptyset$ for every $e_i\in G$. Let's remove $e_j'$ from $H$ and also remove every vertex in $e_j'$ from $G$. Then $e_j\neq\emptyset$, so $e_j$ still contains at least one vertex $v$. This similarly leads to a contradiction, as $e_j'=e_j$ violates the Sperner property of $H$. Therefore, no such $e_j$ can exist.
\end{proof}

\begin{algorithm}[p]
\caption{Algorithm \FKA{} for \TransHyp{} (cf. \cite{FK96})}\label{alg:fka}
\begin{algorithmic}[1]
\State \textbf{Input:} Hypergraphs $G=(V(G),E(G))$ and $H=(V(H),E(H))$
\State \textbf{Output:}
\If{$G=\Tr(H)$}
\State \textbf{TRUE}
\Else
\State Return a witness $e_k\in E(G)$ such that $\exists e_l\in E(H)$ with $e_k\cap e_l=\emptyset$
\EndIf
\Statex \hfill$\triangleright$ Step 1: Sperner reduction
\State Sperner reduce $G$ and $H$ (delete all non-Sperner edges).
\Statex \hfill$\triangleright$ Step 2: Check conditions ($i$)-($iv$)
\State Check the following conditions:
\Statex \hspace{\algorithmicindent}\begin{tabular}{rl}
(i) & $V(G)=V(H)=V$\\
(ii) & $\max\{|e_i|:e_i\in E(G)\}\leq|H|$ and $\max\{|e_j|:e_j\in E(H)\}\leq|G|$\\
(iii) & $e_i\cap e_j\neq\emptyset,\forall e_i\in E(G),e_j\in E(H)$\\
(iv) & $\sum_{e_i\in E(G)}2^{-|e_i|}+\sum_{e_j\in E(H)}2^{-|e_j|}\geq 1$
\end{tabular}
\If{any condition ($i$)-($iv$) fails}
\State $G\neq\Tr(H)$
\State Return $e_k$ as a witness in polynomial time.
\EndIf
\Statex \hfill$\triangleright$ Step 3: Handle trivial cases
\If{$|G|\cdot|H|\leq 1$}
\State Return $e_k$ by a trivial check in $\mathcal{O}(1)$ time.
\EndIf
\Statex \hfill$\triangleright$ Step 4: Recursive decomposition
\If{$|G|\cdot|H|\geq 2$}
\State Select a vertex $v_i\in V$ such that $\epsilon(v_i)\geq\frac{1}{\log(|G|+|H|)}$
\State Define:
\Statex \hspace{\algorithmicindent}\begin{minipage}{0.8\linewidth}
\begin{itemize}
\item $G_0=\{e_j-\{v_i\}\mid v_i\in e_j,e_j\in G\}$
\item $G_1=\{e_j\mid v_i\notin e_j,e_j\in G\}$
\item $H_0=\{e_k-\{v_i\}\mid v_i\in e_k,e_k\in H\}$
\item $H_1=\{e_k\mid v_i\notin e_k,e_k\in H\}$
\end{itemize}
\end{minipage}
\State Call \FKA{} on the following subproblems:
\State \textbf{L}: $(G_1,H_0\vee H_1)$
\State \textbf{R}: $(H_1,G_0\vee G_1)$
\If{\textbf{L} AND \textbf{R} return \textbf{TRUE}}
\State Return \textbf{TRUE} ($G$ and $H$ are transversal hypergraphs).
\Else
\State Return witness $e_k$ in polynomial time.
\EndIf
\EndIf
\end{algorithmic}
\end{algorithm}

\begin{proposition}\label{prop:sum}
\[
\sum_{e_i\in E(G)}2^{-|e_i|}+\sum_{e_j\in E(H)}2^{-|e_j|}\geq 1
\]
\end{proposition}

\begin{proof}
Let $G$ and $H$ be Sperner hypergraphs that are transversal to each other. This implies that each minimal transversal of $G$ corresponds to an edge in $H$, and each minimal transversal of $H$ corresponds to an edge in $G$. Suppose that the sum is less than 1. This would imply that there exists an edge $e_i\in E(G)$ such that
\[
\forall v_k\in e_i,\exists e_j\in H:v_k\notin e_j.
\]
In this case, some vertices in $e_i$ would not intersect any edge in $H$, which contradicts the assumption that $H$ is a transversal of $G$.
\end{proof}

\subsection{Recursive Decomposition}

Let $\epsilon\in(0,1]$. A vertex $v_i\in\{v_1,\dots,v_n\in V\}$ is said to have frequency at least $\epsilon$ in hypergraphs $G$ and $H$ if:
\[
\epsilon(v_i)\leq\max\left\{\frac{\#\{e_j\in E(G):v_i\in e_j\}}{|E(G)|},\frac{\#\{e_j\in E(H):v_i\in e_j\}}{|E(H)|}\right\}.
\]

We denote that a vertex $v_i$ has frequency $\geq\epsilon$ in $G$ and $H$ if $|E(G)|\cdot|E(H)|\geq 1$ and $v_i$ has a frequency of at least $\epsilon$ in either $G$ or $H$.

This frequency condition ensures the existence of a vertex $v_i$ with these properties, referred to as a \emph{frequent vertex}. For $G$ and $H$, the vertex with the highest frequency in either hypergraph is called the \emph{most frequent vertex} (MFV).

As described in \ref{alg:fka}, for a most frequent vertex $v\in V$, the hypergraphs $G$ and $H$ can be rewritten as:
\[
G=(\{v\}\wedge G_0)\vee G_1,\quad H=(\{v\}\wedge H_0)\vee H_1
\]
where:
\begin{gather*}
G_0=\{e_j-\{v\}\mid v\in e_j,e_j\in G\},\quad G_1=\{e_j\mid v\notin e_j,e_j\in G\}\\
H_0=\{e_k-\{v\}\mid v\in e_k,e_k\in H\},\quad H_1=\{e_k\mid v\notin e_k,e_k\in H\}
\end{gather*}

The algorithm then reduces the problem to the following smaller subproblems:
\[
\textbf{L:}\ \ (G_1,H_0\wedge H_1),\qquad\textbf{R:}\ \ (H_1,G_0\wedge G_1)
\]

If both subproblems \textbf{L} and \textbf{R} return TRUE, the algorithm confirms that $G$ and $H$ are transversal hypergraphs. Otherwise, the algorithm identifies a hyperedge $w$ in polynomial time that violates one of the transversal properties.

We define the subhypergraphs as input for the subproblems generated by \FKA{} recursion as the following:

\begin{definition}\label{def:contraction}
Let $H=(V,E)$ be a hypergraph and $v\in V$ be a vertex. We define the \emph{contraction subhypergraph} of $v$ as the induced subhypergraph $H/v=(V_v,E_v)$ where $V_v=V\setminus\{v\}$, and $E_v=\{e'\subseteq e\in E\mid e'=e\cap V_v,\forall e\in E\}$. For hypergraph $G$, this is equivalent to the $G_0\wedge G_1$.
\end{definition}

\begin{definition}\label{def:deletion}
Let $H=(V,E)$ be a hypergraph and $v\in V$ be a vertex. The \emph{deletion subhypergraph} of $v$ is defined as the subhypergraph $H\setminus v=(V_v,E_v)$ where $V_v=V\setminus\{v\}$, and $E_v=\{e\in E\mid v\notin e\}$.
\end{definition}

\begin{example}
Consider the hypergraph $H=(V,E)$ where $V=\{v_1,v_2,v_3,v_4\}$ and $E=\{\{v_1,v_2\},\{v_2,v_3,v_4\},\{v_1,v_3\}\}$.

Let $V'=V\setminus\{v_1\}$ be a subset of vertices. The contraction $H/v=(V',E/v_1)$ is given by $V'=V\setminus\{v\}$ and:
\[
E/v_1=\{\{v_2\},\{v_2,v_3,v_4\},\{v_3\}\},
\]
where each edge in $E/v_1$ is the intersection of an edge in $E$ with $V'$. The deletion $H\setminus v_1=(V',E\setminus v_1)$ is given by :
\[
E\setminus v_1=\{\{v_2,v_3,v_4\}\},
\]
where $E\setminus v_1$ consists of edges from $E$ that do not contain $v_1$.
\end{example}

Suppose, $H=(V,E)$ is a Sperner hypergraph, where $v\in V$. $H/v$ is the $v$ induced contraction subhypergraph. $H/v$ does not necessarily retain the Sperner property. However, the $v$ induced deletion subhypergraph $H\setminus v$ is just a set of edges not containing $v$. Thus, their Sperner property is retained.

\section{\FKA{} Bound}

\subsubsection*{Bound on Hyperedge Size}

In the base operation of \FKA, Proposition \ref{prop:sum} ensures that transversal hypergraphs contain hyperedges of logarithmic length (in terms of the number of vertices). Using this property, the size of hyperedges can be bounded by the following result:

\begin{corollary}\label{cor:shortest}
Let $G$ and $H$ be mutually transversal hypergraphs, where $n=|G|+|H|$, and let $s=\min\{|e_l|:e_l\in E(G)\cup E(H)\}$ represent the length of the shortest hyperedge. Then:
\[
s\geq\log n,\quad\text{since }n\cdot 2^{-s}=(|G|+|H|)\cdot 2^{-s}\geq 1.
\]
\end{corollary}
\noindent\textit{\cite{FK96}}

\subsubsection*{Bound on Vertex Frequency}

Both Proposition \ref{prop:sum} and Corollary \ref{cor:shortest} establish bounds on the frequency of the most frequent vertex (MFV). This is formalized as follows:

\begin{lemma}\label{lem:frequency}
For transversal hypergraphs $G$ and $H$, with $|G|\cdot|H|\geq 1$ and $n=|G|+|H|$, there exists a vertex $v_i\in V(G)=V(H)$ with frequency:
\[
\epsilon(v_i)\geq\frac{1}{\log n}.
\]
\end{lemma}
\noindent\textit{\cite{FK96}}

These bounds provide a lower limit on the performance of \FKA. Identifying better upper bounds is an open challenge, particularly for hypergraph classes where \FKA{} struggles.

\begin{observation}\label{obs:tight}
The bounds in Corollary \ref{cor:shortest} and Lemma \ref{lem:frequency} are tight up to a factor of 2. There exist transversal hypergraph pairs with arbitrarily large $n$ such that:
\[
s=\min\{|e_l|:e_l\in E(G)\cup E(H)\}\geq\frac{\log n}{2},\quad\text{and}
\]
\[
\max\left(\epsilon_i^G,\epsilon_i^H\right)<\frac{2}{\log n},\quad\text{for any }i=1,\dots,n.
\]
\end{observation}
\noindent\textit{\cite{FK96}}

\subsubsection*{Bound on Number of Subproblems}

In each recursive step, \FKA{} splits the problem into two smaller subproblems, denoted \textbf{L} and \textbf{R}. The algorithm's running time is bounded (up to a factor of $n^{\mathcal{O}(1)}$) by the total number of recursive calls, denoted $\alpha(G,H)$. The following lemma establishes this bound:

\begin{lemma}\label{lem:subproblems}
For hypergraphs $G$ and $H$:
\[
\alpha(G,H)\leq n^{4\log^2 n}.
\]
\end{lemma}
\noindent\textit{\cite{FK96}}

These results lead to the first quasi-polynomial bound for solving \TransHyp{} using \FKA, as described in \cite{FK96}:

\begin{theorem}
\textit{\texttt{Trans-Hyp}} for $G$ and $H$ can be solved in:
\[
n^{4\cdot\chi(n)+\mathcal{O}(1)}\text{ time, where }\chi(n)^{\chi(n)}=n\text{ and }\chi(n)\sim\log\log n=\mathcal{O}(\log n).
\]
\end{theorem}

This result implies that solving \TransHyp{} is possible within $\mathcal{O}(\log n)$ time, which corresponds to the bound for FK-B. Additionally, given a hypergraph $G$ and a proper subhypergraph of its transversal hypergraph $H$, denoted $H^c\subseteq H$, a new hyperedge $e_i\in H\setminus H^c$ can be generated in incremental time:
\[
n^{4\cdot\chi(n)+\mathcal{O}(1)}.
\]
